% ── chemistry-handout-preamble.tex ── 化学竞赛讲义 LaTeX 导言区（v7 · 精排版）
\documentclass[a4paper,11pt]{ctexart}

% ── 字体优化 ──
\setCJKmainfont[AutoFakeBold=true,BoldFont=SimHei]{SimSun}
\setCJKsansfont{SimHei}

% 化学式宏（mhchem）
\usepackage[version=4]{mhchem}

% ── 数学 ──
\usepackage{amsmath,amssymb}

% ── 图片 ──
\usepackage{graphicx}

% ── 表格（跨页+自动列宽）──
\usepackage{booktabs}
\usepackage{xltabular}
\usepackage{array}
\newcolumntype{L}{>{\raggedright\arraybackslash}X}
\newcolumntype{C}{>{\centering\arraybackslash}X}
\newcolumntype{R}{>{\raggedleft\arraybackslash}X}

% ── 浮动体（[H] 强制定位）──
\usepackage{float}

% ── 页面 ──
\usepackage[includeheadfoot,margin=2.2cm,headheight=14pt]{geometry}
\setlength{\textwidth}{16.6cm}

% ── 行间距 ──
\usepackage{setspace}
\setstretch{1.15}
\setlength{\parskip}{0.3em}
\setlength{\parindent}{2em}

% ── 防止表格跨页 ──
\usepackage{longtable}
\setlength{\LTpre}{0pt}
\setlength{\LTpost}{0pt}

% ── 防止图片跨页（尽量）──
\renewcommand{\textfraction}{0.1}
\renewcommand{\topfraction}{0.9}
\renewcommand{\bottomfraction}{0.9}
\setcounter{totalnumber}{10}
\setcounter{topnumber}{10}
\setcounter{bottomnumber}{10}

% ── 超链接(无红框) ──
\usepackage[colorlinks=true,linkcolor=black,citecolor=black,urlcolor=black]{hyperref}

% ── 颜色 ──
\usepackage[dvipsnames,svgnames]{xcolor}
\definecolor{chapterblue}{RGB}{23,50,77}
\definecolor{insightbg}{RGB}{235,240,250}
\definecolor{warningbg}{RGB}{255,245,235}
\definecolor{tipbg}{RGB}{235,250,235}

% ── 页眉页脚 ──
\usepackage{fancyhdr}
\pagestyle{fancy}
\fancyhf{}
\fancyhead[L]{\small\color{chapterblue}\leftmark}
\fancyhead[R]{\small\color{chapterblue}\thepage}
\renewcommand{\headrule}{{\color{chapterblue}\hrule height 0.8pt}}

% ── 章节编号与样式 ──
\setcounter{secnumdepth}{1}
\ctexset{
  section = {
    format = \Large\bfseries\heiti\color{chapterblue},
    aftername = \hspace{0.5em},
    beforeskip = 2ex,
    afterskip = 1ex,
  },
  subsection = {
    format = \large\bfseries\heiti,
    beforeskip = 1.5ex,
    afterskip = 0.5ex,
  },
  subsubsection = {
    format = \normalsize\bfseries\heiti,
    beforeskip = 1ex,
    afterskip = 0.3ex,
  },
  paragraph = {
    format = \small\bfseries,
    beforeskip = 0.5ex,
    afterskip = 0.2ex,
  },
}

% ── blockquote（引用/提示框）缩减 ──
\usepackage{etoolbox}
\AtBeginEnvironment{quote}{\small\setstretch{1.05}\leftskip=1em\rightskip=1em}

% ── 教学提示框（浅底色版）──
\newenvironment{insightbox}{%
  \par\smallskip
  \noindent{\small\bfseries\color{RoyalBlue}🧠 认知冲突}\par
  \leavevmode\colorbox{insightbg}{\parbox{\dimexpr\textwidth-2\fboxsep\relax}{%
  \ignorespaces}}%
}{\par\smallskip}

\newenvironment{warningbox}{%
  \par\smallskip
  \noindent{\small\bfseries\color{BrickRed}⚠️ 易错点}\par
  \leavevmode\colorbox{warningbg}{\parbox{\dimexpr\textwidth-2\fboxsep\relax}{%
  \ignorespaces}}%
}{\par\smallskip}

\newenvironment{tipbox}{%
  \par\smallskip
  \noindent{\small\bfseries\color{ForestGreen}💡 理解提示}\par
  \leavevmode\colorbox{tipbg}{\parbox{\dimexpr\textwidth-2\fboxsep\relax}{%
  \ignorespaces}}%
}{\par\smallskip}

% ── 列表（答案区加间距）──
\usepackage{enumitem}
\setlist{nosep}
% 专门用于参考答案的列表（每项之间空一行）
\newlist{answerlist}{enumerate}{1}
\setlist[answerlist]{label=\arabic*., leftmargin=2em, itemsep=0.8em}

% ── 防止内容溢出 ──
\usepackage{microtype}
\setlength{\emergencystretch}{2em}
\setlength{\hfuzz}{2pt}
\setlength{\overfullrule}{0pt}
\sloppy

% ── pandoc 兼容 ──
\providecommand{\tightlist}{\setlength{\itemsep}{0pt}\setlength{\parskip}{0pt}}
\newenvironment{Shaded}{}{}
\newenvironment{Highlighting}{}{}
\newcommand{\NormalTok}[1]{#1}
\newcommand{\CommentTok}[1]{#1}
\newcommand{\KeywordTok}[1]{#1}
\newcommand{\StringTok}[1]{#1}
\newcommand{\FunctionTok}[1]{#1}
\newcommand{\DataTypeTok}[1]{#1}
\newcommand{\NumberTok}[1]{#1}
\newcommand{\ControlFlowTok}[1]{#1}
\newcommand{\OperatorTok}[1]{#1}
\newcommand{\BaseNTok}[1]{#1}
\newcounter{none}

% ── 图片尺寸限制 ──
\AtBeginDocument{\setkeys{Gin}{width=\textwidth,keepaspectratio}}

% ── 防止图片跨节 ──
\usepackage{placeins}

% ── 自定义命令 ──
% 练习题标题（去掉括号）
% 例题标题（去掉中括号）
% 这些需要在convert.py预处理中做正则替换
